\documentclass{article}
\usepackage[UTF8]{ctex}
\usepackage{amsmath,amssymb,amsfonts}
\usepackage{enumitem}
\usepackage{geometry}
\geometry{a4paper, margin=1in}
% === 错题本排版优化（不参与题目标识） ===
\setlength{\parindent}{0pt}
\setlength{\parskip}{0.35em plus 0.10em minus 0.08em}
\setlength{\emergencystretch}{3em}
\setlist{itemsep=0.25em, topsep=0.40em, parsep=0pt}
\allowdisplaybreaks[2]
% === 错题本排版优化结束 ===
\begin{document}
\section*{广西大学2024年高等代数真题}
\section*{152．广西大学2024年高等代数试卷}
一．填空题．每题6分，共48分．\\

\subsection{题目1}
1．若 \(x^{2}+1 \mid x^{4}+k x^{2}+1\) ，则 \(k=\) \(\_\_\_\_\)．

\subsection{题目2}
2．设向量组

\[
\begin{gathered}
\alpha_{1}=(1,0,3,4,3)^{\prime}, \alpha_{2}=(3,-1,2,1,3)^{\prime}, \alpha_{3}=(-1,1,0,5,2)^{\prime} \\
\alpha_{4}=(3,0,5,10,8)^{\prime}, \alpha_{5}=(-1,0,1,-2,-2)^{\prime}
\end{gathered}
\]

则 \(\alpha_{1}, \alpha_{2}, \alpha_{3}, \alpha_{4}, \alpha_{5}\) 的秩是 \(\_\_\_\_\)．

\subsection{题目3}
3．已知二次型 \(x_{1}^{2}+x_{2}^{2}+5 x_{3}^{2}+2 t x_{1} x_{2}-2 x_{1} x_{3}+4 x_{2} x_{3}\) 是正定的，则 \(t\) 满足 \(\_\_\_\_\)．

\subsection{题目4}
4．设 \(P\) 为数域，在 \(P^{4}\) 中，令

\[
\begin{gathered}
W_{1}=\left\{\left(x_{1}, x_{2}, x_{3}, x_{4}\right) \mid x_{1}-2 x_{2}+2 x_{4}=0, x_{1}+2 x_{3}=0\right\} \\
W_{2}=\left\{\left(x_{1}, x_{2}, x_{3}, x_{4}\right) \mid x_{1}-4 x_{2}-2 x_{3}+4 x_{4}=0 .\right\}
\end{gathered}
\]

则 \(W_{1} \cap W_{2}\) 的维数是 \(\_\_\_\_\)．

\subsection{题目5}
5．若 \(A B=B A\) ，矩阵 \(B\) 就称与 \(A\) 可交换，则所有与 \(\left(\begin{array}{lll}0 & 1 & 0 \\ 0 & 0 & 1 \\ 0 & 0 & 0\end{array}\right)\) 可交换的矩阵为 \(\_\_\_\_\)．\\
\(\_\_\_\_\)

\subsection{题目6}
6．若方程组 \(\left\{\begin{array}{l}x_{1}+x_{2}+x_{3}=3 ; \\ x_{1}+2 x_{2}-a x_{3}=9 ; \\ 2 x_{1}-x_{2}+3 x_{3}=6 .\end{array}\right.\) 有唯一解，则 \(a\) 满足的条件是 \(\_\_\_\_\)

\subsection{题目7}
7．设 \(P\) 为数域，在 \(P^{3}\) 中给出一组基 \(\alpha_{1}=(-1,0,2), \alpha_{2}=(0,1,1), \alpha_{3}=(3,-1,0)\) ，定义线性变换 \(\sigma\)如下：\(\sigma\left(\alpha_{1}\right)=(-5,0,3), \sigma\left(\alpha_{2}\right)=(0,-1,6), \sigma\left(\alpha_{3}\right)=(-5,-1,9)\) 。则 \(\sigma\) 在基 \(\alpha_{1}, \alpha_{2}, \alpha_{3}\) 下的矩阵为\\
\(\_\_\_\_\)．

\subsection{题目8}
8．设二阶矩阵 \(\left(\begin{array}{cc}-2 & 1 \\ 0 & t\end{array}\right)\) 与 \(\left(\begin{array}{cc}-10 & -4 \\ 26 & 11\end{array}\right)\) 相似，则 \(t=\)

二．（12 分）计算 \(n\) 阶行列式

\[
D_{n}=\left|\begin{array}{ccccc}
x & a & a & \cdots & a \\
-a & x & a & \cdots & a \\
-a & -a & x & \cdots & a \\
\vdots & \vdots & \vdots & & \vdots \\
-a & -a & -a & \cdots & x
\end{array}\right| .
\]

三．（18 分）设 \(V_{1}, V_{2}, \cdots, V_{s}\) 是数域 \(P\) 上线性空间 \(V\) 的 \(s\) 个非平凡子空间，证明：\(V\) 中至少存在向量\(\alpha\) ，使得 \(\alpha \notin V_{i}, i=1,2, \cdots, s\) ．

四．（12 分）设 \(\sigma\) 是数域 \(P\) 上线性空间 \(V\) 的线性变换，且 \(\sigma^{2}=\varepsilon\)（ \(\varepsilon\) 是 \(V\) 的恒等变换）。证明：对 \(V\)中每个向量 \(\alpha\) ，存在 \(\alpha_{1}, \alpha_{2} \in V\) ，使得 \(\sigma\left(\alpha_{1}\right)=\alpha_{1}, \sigma\left(\alpha_{2}\right)=-\alpha_{2}\) ，且 \(\alpha\) 可唯一地表示成 \(\alpha_{1}\) 与 \(\alpha_{2}\) 之和（即 \(\alpha=\alpha_{1}+\alpha_{2}\) 表示唯一）。

五．（12 分）设有齐次线性方程组

\[
\left\{\begin{array}{c}
a_{11} x_{1}+a_{12} x_{2}+\cdots+a_{1 n} x_{n}=0 ; \\
a_{21} x_{1}+a_{22} x_{2}+\cdots+a_{2 n} x_{n}=0 ; \\
\cdots \cdots \\
a_{n-1,1} x_{1}+a_{n-1,2} x_{2}+\cdots+a_{n-1, n} x_{n}=0 .
\end{array}\right.
\]

\(M_{i}(i=1,2, \cdots, n)\) 为系数矩阵 \(A\) 划去地 \(i\) 列剩下的 \((n-1) \times(n-1)\) 矩阵的行列式。证明：如果秩\((A)=n-1\) ，则 \(\eta_{0}=\left(M_{1},-M_{2}, \cdots,(-1)^{n-1} M_{n}\right)\) 是方程组的一个基础解系。

六．（18 分）设 \(A=\left(a_{i j}\right)_{n \times n}, B=\left(b_{i j}\right)_{n \times n}, C=\left(a_{i j} b_{i j}\right)_{n \times n}\) ．证明：若 \(A, B\) 均是正定矩阵，则 \(C\) 也是正定矩阵。

七．（18 分）设 \(m\) 为正整数，\(f(x)\) 和 \(g(x)\) 为数域 \(P\) 上的非零多项式，证明：\(g^{m}(x) \mid f^{m}(x)\) 的充分必要条件是 \(g(x) \mid f(x)\) ．

八．（12 分）设 \(\alpha\) 是欧氏空间 \(V\) 的一个非零向量，\(\alpha_{1}, \alpha_{2}, \cdots, \alpha_{m} \in V\) 满足

\[
\left(\alpha_{i}, \alpha\right)>0(i=1,2, \cdots, n) ;\left(\alpha_{i}, \alpha_{j}\right) \leq 0(i, j=1,2, \cdots, m ; i \neq j) .
\]

其中符号 \((x, y)\) 表示向量 \(x, y\) 的内积．证明：\(\alpha_{1}, \alpha_{2}, \cdots, \alpha_{m}\) 线性无关．
\end{document}
